% GENERATED FILE. Edit canonical lessons, then run npm run generate:books.
% canonical-source-hash: fnv1a64:54d1cc17a706da04
% canonical-lessons: MR-C197-savay-hone, MR-C197-ghai-karne, MR-C197-javal-yene, MR-C197-gappa-marne, MR-C197-pocavne

\chapter{To get used to, To hurry, To come near, To chat, To deliver}
\label{ch:mr-to-get-used-to-to-hurry-to-come-near-to-chat-to-deliver}
\hlchaptermodality{197}

\begin{chapteropening}
\textbf{By the end of this chapter:} \emph{I can say to get used to, to hurry, to come near, to chat and to deliver.}

Complete the last lesson of chapter 197: I can say to get used to, to hurry, to come near, to chat and to deliver.
\end{chapteropening}

\section[\texorpdfstring{\mr{सवय} \mr{होणे} (savay hoṇe)}{सवय होणे (savay hoṇe)}]{\mr{सवय} \mr{होणे} (savay hoṇe) — to get used to}

\label{lesson:MR-C197-savay-hone}



\begin{quote}
Before the new one: say the Marathi for to introduce, then the Marathi for to lend.
\end{quote}

\subsection*{You'll want to know: \mr{सवय} \mr{होणे}}
\textbf{\mr{सवय} \mr{होणे}} — \emph{savay hoṇe} — \textquotedblleft{}to get used to\textquotedblright{}.

\begin{grammarlens}[title={What you've built}]
One more word for this chapter.
\end{grammarlens}

\subsection*{Your turn}
\begin{itemize}
\raggedright
  \item \emph{Say it:} \emph{savay hoṇe}
  \item \emph{Say it:} \emph{savay hoṇe}, once more
  \item \emph{Say it:} say \emph{savay hoṇe}
  \item \emph{Recall:} say the Marathi for to introduce, then the Marathi for to lend, then say \emph{savay hoṇe} again
\end{itemize}

\begin{tcolorbox}[breakable,colback=teal!4,colframe=teal!35!black,title={Before you move on}]
What does \mr{सवय} \mr{होणे} mean? (\textquotedblleft{}To get used to\textquotedblright{}.) Say it once more.
\end{tcolorbox}
\section[\texorpdfstring{\mr{घाई} \mr{करणे} (ghāī karṇe)}{घाई करणे (ghāī karṇe)}]{\mr{घाई} \mr{करणे} (ghāī karṇe) — to hurry}

\label{lesson:MR-C197-ghai-karne}



\begin{quote}
Before the new one: say the Marathi for to lend, then the Marathi for to get used to.
\end{quote}

\subsection*{You'll want to know: \mr{घाई} \mr{करणे}}
\textbf{\mr{घाई} \mr{करणे}} — \emph{ghāī karṇe} — \textquotedblleft{}to hurry\textquotedblright{}.

\begin{grammarlens}[title={What you've built}]
One more word for this chapter.
\end{grammarlens}

\subsection*{Your turn}
\begin{itemize}
\raggedright
  \item \emph{Say it:} \emph{ghāī karṇe}
  \item \emph{Say it:} \emph{ghāī karṇe}, once more
  \item \emph{Say it:} say \emph{ghāī karṇe}
  \item \emph{Recall:} say the Marathi for to lend, then the Marathi for to get used to, then say \emph{ghāī karṇe} again
\end{itemize}

\begin{tcolorbox}[breakable,colback=teal!4,colframe=teal!35!black,title={Before you move on}]
What does \mr{घाई} \mr{करणे} mean? (\textquotedblleft{}To hurry\textquotedblright{}.) Say it once more.
\end{tcolorbox}
\section[\texorpdfstring{\mr{जवळ} \mr{येणे} (javaḷ yeṇe)}{जवळ येणे (javaḷ yeṇe)}]{\mr{जवळ} \mr{येणे} (javaḷ yeṇe) — to come near}

\label{lesson:MR-C197-javal-yene}



\begin{quote}
Before the new one: say the Marathi for to get used to, then the Marathi for to hurry.
\end{quote}

\subsection*{You'll want to know: \mr{जवळ} \mr{येणे}}
\textbf{\mr{जवळ} \mr{येणे}} — \emph{javaḷ yeṇe} — \textquotedblleft{}to come near\textquotedblright{}.

\begin{grammarlens}[title={What you've built}]
One more word for this chapter.
\end{grammarlens}

\subsection*{Your turn}
\begin{itemize}
\raggedright
  \item \emph{Say it:} \emph{javaḷ yeṇe}
  \item \emph{Say it:} \emph{javaḷ yeṇe}, once more
  \item \emph{Say it:} say \emph{javaḷ yeṇe}
  \item \emph{Recall:} say the Marathi for to get used to, then the Marathi for to hurry, then say \emph{javaḷ yeṇe} again
\end{itemize}

\begin{tcolorbox}[breakable,colback=teal!4,colframe=teal!35!black,title={Before you move on}]
What does \mr{जवळ} \mr{येणे} mean? (\textquotedblleft{}To come near\textquotedblright{}.) Say it once more.
\end{tcolorbox}
\section[\texorpdfstring{\mr{गप्पा} \mr{मारणे} (gappā mārṇe)}{गप्पा मारणे (gappā mārṇe)}]{\mr{गप्पा} \mr{मारणे} (gappā mārṇe) — to chat}

\label{lesson:MR-C197-gappa-marne}



\begin{quote}
Before the new one: say the Marathi for to hurry, then the Marathi for to come near.
\end{quote}

\subsection*{You'll want to know: \mr{गप्पा} \mr{मारणे}}
\textbf{\mr{गप्पा} \mr{मारणे}} — \emph{gappā mārṇe} — \textquotedblleft{}to chat\textquotedblright{}.

\begin{grammarlens}[title={What you've built}]
One more word for this chapter.
\end{grammarlens}

\subsection*{Your turn}
\begin{itemize}
\raggedright
  \item \emph{Say it:} \emph{gappā mārṇe}
  \item \emph{Say it:} \emph{gappā mārṇe}, once more
  \item \emph{Say it:} say \emph{gappā mārṇe}
  \item \emph{Recall:} say the Marathi for to hurry, then the Marathi for to come near, then say \emph{gappā mārṇe} again
\end{itemize}

\begin{tcolorbox}[breakable,colback=teal!4,colframe=teal!35!black,title={Before you move on}]
What does \mr{गप्पा} \mr{मारणे} mean? (\textquotedblleft{}To chat\textquotedblright{}.) Say it once more.
\end{tcolorbox}
\section[\texorpdfstring{\mr{पोचवणे} (pocavṇe)}{पोचवणे (pocavṇe)}]{\mr{पोचवणे} (pocavṇe) — to deliver}

\label{lesson:MR-C197-pocavne}



\begin{quote}
Before the new one: say the Marathi for to come near, then the Marathi for to chat.
\end{quote}

\subsection*{You'll want to know: \mr{पोचवणे}}
\textbf{\mr{पोचवणे}} — \emph{pocavṇe} — \textquotedblleft{}to deliver\textquotedblright{}.

\begin{grammarlens}[title={What you've built}]
One more word for this chapter.
\end{grammarlens}

\subsection*{Your turn}
\begin{itemize}
\raggedright
  \item \emph{Say it:} \emph{pocavṇe}
  \item \emph{Say it:} \emph{pocavṇe}, once more
  \item \emph{Say it:} say \emph{pocavṇe}
  \item \emph{Recall:} say the Marathi for to come near, then the Marathi for to chat, then say \emph{pocavṇe} again
\end{itemize}

\begin{tcolorbox}[breakable,colback=teal!4,colframe=teal!35!black,title={Before you move on}]
What does \mr{पोचवणे} mean? (\textquotedblleft{}To deliver\textquotedblright{}.) Say it once more.
\end{tcolorbox}
